% ECONOMICS_PROBLEM: {"id": "C78-511", "title": "The proposal-total spectrum of doubly prime instances", "jel_code": "C78", "source_id": "OP-0206"}
% !TeX program = xelatex
\documentclass[11pt,a4paper]{article}
\usepackage[margin=23mm]{geometry}
\usepackage{fontspec}
\setmainfont{Latin Modern Roman}
\usepackage{amsmath,amssymb,mathtools}
\usepackage{xcolor,enumitem,booktabs,longtable,array,xurl,hyperref}
\definecolor{muted}{HTML}{53636C}
\hypersetup{colorlinks=true,urlcolor=blue,linkcolor=blue}
\setlength{\parindent}{0pt}
\setlength{\parskip}{5pt}
\newcommand{\R}{\mathbb R}
\newcommand{\E}{\mathbb E}
\newcommand{\N}{\mathbb N}
\newcommand{\ind}{\mathbf 1}
\newcommand{\Var}{\operatorname{Var}}
\newcommand{\Cov}{\operatorname{Cov}}
\newcommand{\supp}{\operatorname{supp}}
\DeclareMathOperator*{\argmin}{arg\,min}
\DeclareMathOperator*{\argmax}{arg\,max}
\newcommand{\entrymeta}[3]{{\sffamily\small\color{muted}\textbf{\texttt{#1}}\quad|\quad #2\par #3\par}}
\newcommand{\Problemref}[1]{\texttt{#1}}
\newcommand{\JELref}[2]{\texttt{#2} (JEL \texttt{#1})}
\begin{document}
\section*{The proposal-total spectrum of doubly prime instances}
\textbf{Problem ID:} \texttt{C78-511}\quad \textbf{JEL:} \texttt{C78}\par
% BEGIN ECONOMICS STATEMENT
\label{problem:OP-0206}
% GAME_THEORY_PROBLEM: {"local_id": "GT-076", "original_number": 76, "kind": "P", "entry_kind": "statement_only", "published": false, "review_required": true, "category": "game_theory"}
\label{game:GT-076}
{\sffamily\small\color{muted}
\noindent\textbf{ID: GT-076.}\quad\textbf{Category:} Deferred acceptance combinatorics.
\quad\textbf{Kind: P.}\quad\textbf{Status:} Conjecture in the cited source.
\par}
\subsection*{Definitions and mathematical statement}
Let \(I\) be a strict complete stable-marriage instance on \(n\)
proposers \(P\) and \(n\) receivers \(R\). Serial deferred acceptance
starts with all proposers unheld. At each legal step an unheld proposer
makes its next untried proposal, and that receiver keeps the preferred
of its previous holder and the new proposer.
A state \((h,q)\) records each receiver's holder
\(h:R\to P\cup\{\bot\}\) and each proposer's proposal counter
\(q:P\to\{0,\ldots,n\}\). Let \(\operatorname{Reach}(I)\) be all states
attainable from the empty initial state by legal steps.
The terminal counter vector \(q^\ast\) is independent of the legal
serial schedule; define
\[
                         Q^\ast(I)=\sum_{p\in P}q^\ast(p).
\]

For completeness, use these precise primality conditions.
Input-prime means that no partitions
\(P=\bigsqcup_{j=1}^t P_j\), \(R=\bigsqcup_{j=1}^t R_j\),
with \(t\geq2\) and \(|P_j|=|R_j|>0\), make every member of each
paired block rank all opposite-side members of that block above
all outsiders. Define
\[
 F(q)=\{(p,t):p\in P,\ t\in\mathbb Z,\ 0\leq t<q(p)\},\quad
 \mathcal F=\{F(q):(h,q)\in\operatorname{Reach}(I)\},\quad
 E=\bigcup_{X\in\mathcal F}X.
\]
Behaviorally prime means that no partition \(E=E_1\sqcup E_2\)
with both parts nonempty satisfies
\[
 \mathcal F=
 \{X_1\cup X_2:
       X_1\in\{X\cap E_1:X\in\mathcal F\},
       \ X_2\in\{X\cap E_2:X\in\mathcal F\}\}.
\]
Let \(\mathcal D_n\) be the instances satisfying both conditions.

\paragraph{Conjecture.}
For every integer \(n\geq 2\),
\[
 \{Q^\ast(I):I\in\mathcal D_n\}
   =\{2n-1,2n,\ldots,n^2-n+1\}.
\]
Equivalently, every integer between the displayed lower and upper
bounds is attained by some doubly prime instance.

\subsection*{Short English statement}
Does double primality permit every proposal total in the allowed interval?

\subsection*{Variants and scope boundaries}
The total counts proposals in a complete run, including rejected
proposals. It does not depend on choosing a favorable schedule.
The event-family definition is used rather than an experimentally
validated shortcut for testing primality.

\subsection*{Source relationship and status}
This is [S1, Section 7, problem C].
The two endpoints are valid general bounds in this class; the paper
reports attained intermediate values at several small sizes.
Those finite computations do not prove the universal assertion.

\subsection*{References}
\begingroup\footnotesize\raggedright
\noindent[S1] Yoshiteru Ishida,
\emph{Prime Factorisation of Stable-Matching Instances: Uniqueness,
Simultaneous Products, and an Exact Census}, 2026,
arXiv:2610.05617v1.
\url{https://arxiv.org/html/2610.05617v1}.
Locators: Sections 2 and 4; Proposition 21.4;
Section 7, problem C.
\par\endgroup

\subsection*{Common conventions: Source mathematical conventions}

\textbf{Finite games.} For a finite set $A$, $\Delta(A)$ is its probability
simplex. Players choose independent mixed strategies in Nash equilibrium;
correlation is allowed only when explicitly part of the solution concept.
In a game with payoff functions $u_i$ and strategy profile $x$, an additive
$\varepsilon$ Nash equilibrium satisfies
\[
 u_i(x)\geq u_i(x_i',x_{-i})-\varepsilon
 \quad\text{for every player $i$ and every admissible deviation $x_i'$}.
\]
For numerical approximation, payoffs are normalized as locally specified.
An $\varepsilon$ payoff residual is distinct from distance at most
$\varepsilon$ to the set of exact equilibria. Point convergence, convergence
of empirical play, and convergence in distance to an equilibrium set are
also distinct.

\textbf{Computational representations.} Unless stated otherwise, a complexity
question takes rational data with binary encoding; polynomial time is measured
in total bit length. A fixed-accuracy polynomial algorithm is not necessarily
polynomial in $1/\varepsilon$, and neither is necessarily polynomial in
$\log(1/\varepsilon)$. Succinct games and value, demand, or sampling oracles
must declare their interfaces. Exact real solutions require an explicit
representation; an unspecified list of arbitrary real numbers is not a
finite computational output. A claimed algorithmic barrier is meaningful
only for its specified model and reductions.

\textbf{Dynamic games.} Histories, information, observation, and allowable
randomization are part of the model. A stationary strategy depends only on
the current state; a positional strategy in a deterministic graph game
chooses one action at each controlled vertex. Nash equilibrium at an initial
state and equilibrium after every history are different requirements.
An expectation of a pathwise $\liminf$ payoff, a $\liminf$ of expected averages,
a limiting expected average, and a uniform finite-horizon equilibrium are
different criteria. None is silently substituted for another.

\textbf{Allocation and mechanisms.} A complete allocation assigns every item
exactly once unless otherwise stated. Zero-valued goods, negative values,
capacity constraints, feasibility, lotteries, and copies of goods affect
fairness definitions and are specified locally. Dominant-strategy incentive
compatibility, Bayesian incentive compatibility, universal truthfulness,
and truthfulness in expectation impose different inequalities. Whenever
an objective ratio has zero denominator, its convention is stated or those
instances are excluded.

\textbf{Infinite-dimensional models.} Expectations, integrals, stochastic
equations, and optimization objectives require their local regularity,
integrability, filtrations, and admissible domains. A backward--forward
mean-field system is a boundary-value problem; it is not an initial-value
semiflow without an additional construction. Long-time, large-population,
and vanishing-noise limits require an order of limits and a convergence
topology. If a minimum need not be attained, the objective is its infimum
or supremum and the approximation requirement is stated separately.
% END ECONOMICS STATEMENT
\end{document}
